% Language: English
% Category: REF
% ============================================================================
% APPENDIX BL: REF-DOI
% Digital Object Identifier Reference System
% ============================================================================

\chapter{Digital Object Identifier Reference System}
\label{app:doi-reference}

% ----------------------------------------------------------------------------
% HEADER BLOCK
% ----------------------------------------------------------------------------
\begin{tcolorbox}[colback=gray!5, colframe=gray!50, title=Appendix: BL REF-DOI]
\small
\begin{tabular}{ll}
\textbf{Category:} & REF (Reference Tools) \\
\textbf{Code:} & BL \\
\textbf{Name:} & DOI \\
\textbf{Version:} & 1.0 \\
\textbf{Last Updated:} & 2026-01-27 \\
\textbf{Status:} & ACTIVE \\
\end{tabular}

\vspace{0.5em}
\textbf{Purpose:} Documents the DOI system for EBF bibliography management, including ISO 26324 standard, publisher algorithms, journal-era model, and strategy determination for predictive DOI generation vs. API lookup.

\textbf{Dependencies:} \texttt{journal-database.yaml} (v3.0), \texttt{doi\_strategy.py}, \texttt{bcm\_master.bib}
\end{tcolorbox}

% ----------------------------------------------------------------------------
% CROSS-REFERENCE MAP
% ----------------------------------------------------------------------------
\begin{tcolorbox}[colback=gray!5, colframe=gray!50, title=Cross-Reference Map]
\small
\textbf{Primary Coverage:}
\begin{itemize}[noitemsep]
    \item Bibliography Management (bcm\_master.bib)
    \item Paper SuperKey System (PAP- prefix)
    \item Journal SuperKey System (JRN- prefix)
\end{itemize}

\textbf{Related Appendices:}
\begin{itemize}[noitemsep]
    \item Appendix G (Glossary) -- DOI terminology
    \item Appendix BM (METHOD-API) -- External API integration via GitHub Actions
    \item Appendix BBB (CORE-WHERE) -- Parameter sources from literature
    \item Appendix MS (REF-MODELSPACE) -- Theory catalog with paper references
\end{itemize}

\textbf{Data Sources:}
\begin{itemize}[noitemsep]
    \item \texttt{data/journal-database.yaml} -- Journal patterns and eras
    \item \texttt{scripts/doi\_strategy.py} -- Strategy determination
    \item \texttt{bibliography/bcm\_master.bib} -- 2,324 papers
\end{itemize}
\end{tcolorbox}

% ----------------------------------------------------------------------------
% CHAPTER LINKAGE
% ----------------------------------------------------------------------------
\begin{tcolorbox}[colback=blue!5, colframe=blue!50, title=Chapter Linkage]
\small
\textbf{Primary Function:} This appendix serves as a \emph{reference utility} for all EBF chapters that cite literature. It does not link to a specific chapter but supports the entire framework's bibliography management.

\textbf{Used By:}
\begin{itemize}[noitemsep]
    \item All chapters citing papers from \texttt{bcm\_master.bib}
    \item LIT-Appendices (R, M, O) for author-specific bibliographies
    \item Appendix BBB (CORE-WHERE) for parameter source references
    \item Appendix MS (REF-MODELSPACE) for theory-paper linkages
\end{itemize}

\textbf{Workflow Integration:}
\begin{enumerate}[noitemsep]
    \item New paper entry $\to$ Run \texttt{doi\_strategy.py --key \{key\}}
    \item PATTERN result $\to$ Generate and add DOI
    \item LOOKUP result $\to$ Trigger GitHub Action (see Appendix BM)
    \item Validation $\to$ Verify via \texttt{doi.org}
\end{enumerate}
\end{tcolorbox}

% ----------------------------------------------------------------------------
% SCOPE BOX
% ----------------------------------------------------------------------------
\begin{tcolorbox}[colback=yellow!5, colframe=yellow!50, title=Appendix Scope: What This Appendix Covers]
\small
\textbf{In-Scope:}
\begin{itemize}[noitemsep]
    \item ISO 26324 DOI standard and syntax
    \item Publisher-specific suffix algorithms
    \item Journal-era model for DOI prediction
    \item Strategy determination (PATTERN vs. LOOKUP)
    \item Implementation: \texttt{journal-database.yaml}, \texttt{doi\_strategy.py}
\end{itemize}

\textbf{Out-of-Scope:}
\begin{itemize}[noitemsep]
    \item ISBN lookup for books (see future OI-4)
    \item Automated API integration (see future OI-2)
    \item Citation network analysis
    \item DOI metadata enrichment (DataCite, CrossRef metadata)
\end{itemize}

\textbf{Lieferobjekte (Deliverables):}
\begin{itemize}[noitemsep]
    \item Journal database with 45 journals and DOI eras
    \item Strategy determination algorithm
    \item Per-paper strategy export (CSV)
    \item Predictability rates for prioritization
\end{itemize}
\end{tcolorbox}

% ----------------------------------------------------------------------------
% ABSTRACT
% ----------------------------------------------------------------------------
\begin{abstract}
This appendix documents the Digital Object Identifier (DOI) system as implemented in the Evidence-Based Framework (EBF) bibliography. We present the international standard ISO 26324, analyze publisher-specific suffix algorithms, and develop a journal-era model that enables predictive DOI generation for approximately 24\% of papers. For the remaining 76\%, we document the appropriate lookup strategies (CrossRef, JSTOR, publisher APIs). The journal database covers 45 journals with 2,324 papers, achieving 31\% DOI coverage with a clear path to 100\% through systematic application of the strategy determination algorithm.
\end{abstract}

% ----------------------------------------------------------------------------
% QUICK REFERENCE
% ----------------------------------------------------------------------------
\begin{tcolorbox}[colback=blue!5, colframe=blue!50, title=Quick Reference: DOI Fundamentals]
\small
\textbf{DOI Structure (ISO 26324):}
\begin{verbatim}
doi:10.RRRR/SUFFIX
    |  |    |
    |  |    +-- Suffix (publisher-specific)
    |  +------- Registrant Code (assigned by RA)
    +---------- Directory Indicator (always "10")
\end{verbatim}

\textbf{Example:} \texttt{10.1257/aer.109.4.1256}
\begin{itemize}[noitemsep]
    \item \texttt{10} -- DOI namespace (Handle System)
    \item \texttt{1257} -- American Economic Association
    \item \texttt{aer.109.4.1256} -- Journal.Volume.Issue.Page
\end{itemize}

\textbf{Key Insight:} The prefix is standardized, but the suffix algorithm is publisher-specific and often era-dependent.
\end{tcolorbox}

% ============================================================================
\section{Introduction: The DOI Identification Problem}
\label{sec:doi-intro}
% ============================================================================

The Digital Object Identifier system provides persistent, unique identifiers for scholarly publications. For the EBF bibliography management, DOIs serve three critical functions:

\begin{enumerate}
    \item \textbf{Unique Identification:} Each paper receives a globally unique identifier
    \item \textbf{Persistent Resolution:} DOIs resolve to current URLs even when publishers change
    \item \textbf{Citation Linking:} DOIs enable automated citation networks and impact metrics
\end{enumerate}

\subsection{The Challenge}

While DOI \emph{resolution} is standardized (via \texttt{doi.org}), DOI \emph{construction} is not. Each publisher implements their own suffix algorithm, and these algorithms have evolved over time. This creates a fundamental question:

\begin{tcolorbox}[colback=yellow!10, colframe=orange!50, title=Central Question]
Given a paper's metadata (journal, year, volume, issue, pages), can we \emph{predict} its DOI, or must we \emph{lookup} via external APIs?
\end{tcolorbox}

% ============================================================================
\section{ISO 26324: The International DOI Standard}
\label{sec:iso-standard}
% ============================================================================

\subsection{Standard Evolution}

\begin{table}[h]
\centering
\caption{ISO 26324 Standard Versions}
\begin{tabular}{lll}
\toprule
\textbf{Version} & \textbf{Year} & \textbf{Key Features} \\
\midrule
ISO 26324:2012 & 2012 & Initial standardization \\
ISO 26324:2022 & 2022 & Updated syntax, registration rules \\
ISO 26324:2025 & 2025 & Current version, extended metadata \\
\bottomrule
\end{tabular}
\end{table}

\subsection{DOI Syntax Specification}

The standard defines:

\begin{enumerate}
    \item \textbf{Prefix:} Always \texttt{10.NNNN} where NNNN $\geq$ 1000
    \item \textbf{Separator:} Forward slash (\texttt{/})
    \item \textbf{Suffix:} Any alphanumeric string, case-insensitive
    \item \textbf{Character Set:} \texttt{a-z}, \texttt{A-Z}, \texttt{0-9}, \texttt{-.\_;}()/
    \item \textbf{Length:} No defined limit (optimal: 6-10 characters for suffix)
\end{enumerate}

\subsection{What ISO 26324 Does NOT Standardize}

Critically, the standard does \emph{not} specify:

\begin{itemize}
    \item How suffixes should be constructed
    \item Whether suffixes should encode metadata
    \item How publishers should organize their namespace
\end{itemize}

This gap is the source of the prediction problem.

% ============================================================================
\section{Publisher Suffix Algorithms}
\label{sec:publisher-algorithms}
% ============================================================================

Through empirical analysis of 720 DOIs in the EBF bibliography, we identified six distinct suffix pattern types:

\subsection{Pattern Type Classification}

\begin{table}[h]
\centering
\caption{DOI Suffix Pattern Types}
\label{tab:pattern-types}
\begin{tabular}{llll}
\toprule
\textbf{Type} & \textbf{Pattern} & \textbf{Predictable} & \textbf{Example} \\
\midrule
STRUCTURED & \texttt{j.v.i.p} & Yes & \texttt{10.1257/aer.109.4.1256} \\
ISSN\_YEAR & \texttt{j.issn.yyyy} & Yes & \texttt{10.1111/j.1540-6261.2007.x} \\
NUMERIC & \texttt{NNNNNNN} & No & \texttt{10.2307/1914185} \\
LONG\_CODE & \texttt{15-digit} & No & \texttt{10.1257/000282803322655392} \\
ALPHA\_NUM & \texttt{xxxNNN} & No & \texttt{10.3982/ECTA18709} \\
OTHER & Various & No & \texttt{10.1126/science.aax3100} \\
\bottomrule
\end{tabular}
\end{table}

\subsection{Publisher-Specific Algorithms}

\subsubsection{American Economic Association (10.1257)}

AEA journals show clear era-dependent patterns:

\begin{table}[h]
\centering
\caption{AEA DOI Pattern Evolution}
\begin{tabular}{llll}
\toprule
\textbf{Era} & \textbf{Years} & \textbf{Pattern} & \textbf{Predictable} \\
\midrule
JSTOR Legacy & 1911--1975 & \texttt{10.2307/NNNNNNN} & No \\
AEA Code & 1975--2005 & \texttt{10.1257/NNNNNNNNNNNNNNN} & No \\
Modern & 2005--present & \texttt{10.1257/\{j\}.\{v\}.\{i\}.\{p\}} & Yes \\
\bottomrule
\end{tabular}
\end{table}

\textbf{Modern Pattern:}
\begin{equation}
\text{DOI}_{\text{AER}} = \texttt{10.1257/aer.} \| V \| \texttt{.} \| I \| \texttt{.} \| P_{\text{start}}
\end{equation}

where $V$ is volume, $I$ is issue, and $P_{\text{start}}$ is the starting page.

\subsubsection{University of Chicago Press (10.1086)}

Chicago uses an internal article ID system with \emph{no predictable pattern}:

\begin{equation}
\text{DOI}_{\text{JPE}} = \texttt{10.1086/} \| \text{ArticleID}
\end{equation}

The ArticleID is assigned sequentially by the publisher and cannot be derived from metadata.

\subsubsection{Elsevier (10.1016)}

Elsevier transitioned from PII codes to structured patterns:

\begin{table}[h]
\centering
\caption{Elsevier DOI Pattern Evolution}
\begin{tabular}{llll}
\toprule
\textbf{Era} & \textbf{Years} & \textbf{Pattern} & \textbf{Predictable} \\
\midrule
Old PII & pre-2000 & \texttt{10.1016/ISSN(YY)VVV-S} & No \\
Modern & 2000--present & \texttt{10.1016/j.\{jcode\}.\{yyyy\}.\{mm\}.\{nnn\}} & Partial \\
\bottomrule
\end{tabular}
\end{table}

\subsubsection{Science and Nature (Hash-Based)}

High-impact journals use hash-based suffixes that are \emph{never} predictable:

\begin{align}
\text{DOI}_{\text{Science}} &= \texttt{10.1126/science.} \| \text{Hash} \\
\text{DOI}_{\text{Nature}} &= \texttt{10.1038/} \| \text{Hash}
\end{align}

% ============================================================================
\section{Journal-Era Model}
\label{sec:journal-era-model}
% ============================================================================

\subsection{Model Structure}

We formalize the DOI prediction problem as a journal-era model:

\begin{definition}[Journal-Era Model]
For each journal $j$, define a set of eras $\mathcal{E}_j = \{e_1, e_2, \ldots, e_n\}$ where each era $e_i$ is characterized by:
\begin{itemize}
    \item Year range: $[y_{\text{start}}, y_{\text{end}}]$
    \item DOI prefix: $\pi_i$
    \item Pattern type: $\tau_i \in \{\text{STRUCTURED}, \text{ISSN\_YEAR}, \text{NUMERIC}, \ldots\}$
    \item Predictability: $\rho_i \in \{0, 1\}$
    \item Pattern function: $f_i: \text{Metadata} \to \text{DOI}$ (if $\rho_i = 1$)
\end{itemize}
\end{definition}

\begin{definition}[Predictability Rate]
For journal $j$, the predictability rate $R_j$ is:
\begin{equation}
R_j = \frac{\sum_{e \in \mathcal{E}_j} \rho_e \cdot N_e}{\sum_{e \in \mathcal{E}_j} N_e}
\end{equation}
where $N_e$ is the number of papers in era $e$.
\end{definition}

\subsection{Strategy Determination Algorithm}

\begin{algorithm}[H]
\caption{DOI Strategy Determination}
\label{alg:doi-strategy}
\begin{algorithmic}[1]
\Require Paper $p$ with metadata $(j, y, v, i, \text{pages})$
\Ensure Strategy $\in \{\text{PATTERN}, \text{LOOKUP}\}$ with details
\State $J \gets \textsc{FindJournal}(j)$
\If{$J = \text{null}$}
    \State \Return $(\text{LOOKUP}, \text{CrossRef})$
\EndIf
\State $E \gets \textsc{FindEra}(J, y)$
\If{$E.\rho = 1$} \Comment{Predictable era}
    \If{$\textsc{HasRequiredFields}(p, E.\text{requires})$}
        \State $\text{doi} \gets E.f(p)$ \Comment{Generate DOI}
        \State \Return $(\text{PATTERN}, \text{doi})$
    \EndIf
\EndIf
\State $\text{source} \gets \textsc{DetermineLookupSource}(E.\pi)$
\State \Return $(\text{LOOKUP}, \text{source})$
\end{algorithmic}
\end{algorithm}

\subsection{Lookup Source Determination}

\begin{table}[h]
\centering
\caption{Lookup Source by DOI Prefix}
\begin{tabular}{lll}
\toprule
\textbf{Prefix} & \textbf{Source} & \textbf{Reason} \\
\midrule
\texttt{10.2307} & JSTOR & Historical digitization \\
\texttt{10.1086} & CrossRef/Chicago & Internal codes \\
\texttt{10.1126} & CrossRef/Publisher & Hash-based (Science) \\
\texttt{10.1038} & CrossRef/Publisher & Hash-based (Nature) \\
Other & CrossRef & Default resolution \\
\bottomrule
\end{tabular}
\end{table}

% ============================================================================
\section{Empirical Results}
\label{sec:results}
% ============================================================================

\subsection{Bibliography Statistics}

\begin{table}[h]
\centering
\caption{EBF Bibliography DOI Coverage}
\begin{tabular}{lr}
\toprule
\textbf{Metric} & \textbf{Value} \\
\midrule
Total papers & 2,324 \\
Papers with DOI & 720 (31.0\%) \\
Papers without DOI & 1,604 (69.0\%) \\
Journals in database & 45 \\
\bottomrule
\end{tabular}
\end{table}

\subsection{Strategy Distribution}

For the 1,604 papers without DOI:

\begin{table}[h]
\centering
\caption{DOI Strategy Distribution}
\begin{tabular}{lrrl}
\toprule
\textbf{Strategy} & \textbf{Papers} & \textbf{Percent} & \textbf{Action} \\
\midrule
PATTERN & 173 & 10.8\% & Generate from metadata \\
LOOKUP & 1,152 & 71.8\% & Query external API \\
SKIP & 279 & 17.4\% & Books, not journals \\
\bottomrule
\end{tabular}
\end{table}

\subsection{Predictability by Journal}

\begin{table}[h]
\centering
\caption{Journal Predictability Rates (Top 15)}
\label{tab:predictability}
\begin{tabular}{llrrr}
\toprule
\textbf{SuperKey} & \textbf{Journal} & \textbf{Rate} & \textbf{Error} & \textbf{Type} \\
\midrule
JRN-1016-JPUBE & J. Public Economics & 100\% & 0\% & ISSN\_YEAR \\
JRN-1016-JOEP & J. Economic Psychology & 100\% & 0\% & ISSN\_YEAR \\
JRN-3386-NBER & NBER Working Paper & 100\% & 0\% & ALPHA\_NUM \\
JRN-1162-REST & Rev. Econ. \& Statistics & 100\% & 0\% & STRUCTURED \\
JRN-1257-JEP & J. Economic Perspectives & 96\% & 4\% & STRUCTURED \\
JRN-1257-AER & American Economic Review & 92\% & 8\% & STRUCTURED \\
JRN-1016-JFE & J. Financial Economics & 80\% & 20\% & ISSN\_YEAR \\
JRN-1111-JF & Journal of Finance & 75\% & 25\% & ISSN\_YEAR \\
JRN-1016-GEB & Games \& Econ. Behavior & 75\% & 25\% & ISSN\_YEAR \\
JRN-1257-JEL & J. Economic Literature & 71\% & 29\% & STRUCTURED \\
\midrule
JRN-1093-QJE & Quarterly J. of Economics & 1\% & 99\% & LOOKUP \\
JRN-1086-JPE & J. Political Economy & 0\% & 100\% & LOOKUP \\
JRN-3982-ECTA & Econometrica & 5\% & 95\% & LOOKUP \\
JRN-1126-SCI & Science & 0\% & 100\% & LOOKUP \\
JRN-1038-NAT & Nature & 0\% & 100\% & LOOKUP \\
\bottomrule
\end{tabular}
\end{table}

\subsection{Forecast Error Interpretation}

The \emph{forecast error} represents the proportion of papers for which DOI prediction fails:

\begin{equation}
\text{Forecast Error}_j = 1 - R_j
\end{equation}

\begin{tcolorbox}[colback=green!5, colframe=green!50, title=Practical Interpretation]
\begin{itemize}
    \item \textbf{Error < 10\%:} Use PATTERN strategy; validate sample via CrossRef
    \item \textbf{Error 10--50\%:} Use PATTERN for modern papers; LOOKUP for legacy
    \item \textbf{Error > 50\%:} Use LOOKUP strategy exclusively
\end{itemize}
\end{tcolorbox}

% ============================================================================
\section{Implementation}
\label{sec:implementation}
% ============================================================================

\subsection{Data Files}

\begin{table}[h]
\centering
\caption{DOI System Data Files}
\begin{tabular}{ll}
\toprule
\textbf{File} & \textbf{Purpose} \\
\midrule
\texttt{data/journal-database.yaml} & Journal eras and patterns (v3.0) \\
\texttt{scripts/doi\_strategy.py} & Strategy determination \\
\texttt{scripts/batch\_doi\_by\_journal.py} & Batch DOI generation \\
\texttt{data/doi\_strategy\_export.csv} & Per-paper strategy export \\
\bottomrule
\end{tabular}
\end{table}

\subsection{Usage Examples}

\begin{verbatim}
# Check strategy for a specific paper
python scripts/doi_strategy.py --key fama1993

# Show statistics
python scripts/doi_strategy.py --stats

# Export all strategies to CSV
python scripts/doi_strategy.py --export

# Generate DOIs where possible
python scripts/doi_strategy.py --generate
\end{verbatim}

\subsection{Journal Database Schema}

The \texttt{journal-database.yaml} uses the following structure:

\begin{verbatim}
journals:
  - superkey: "JRN-1257-AER"
    name: "American Economic Review"
    abbreviation: "AER"
    papers_in_ebf: 186
    papers_with_doi: 133
    publisher: "American Economic Association"
    predictability_rate: 0.92
    doi_eras:
      - era: "JSTOR Legacy"
        years: [1911, 1975]
        prefix: "10.2307"
        pattern_type: "NUMERIC"
        predictable: false
      - era: "Modern"
        years: [2005, 2030]
        prefix: "10.1257"
        pattern: "10.1257/aer.{volume}.{issue}.{startpage}"
        pattern_type: "STRUCTURED"
        predictable: true
        requires: [volume, issue, pages]
\end{verbatim}

% ============================================================================
\section{Integration with EBF}
\label{sec:integration}
% ============================================================================

\subsection{SuperKey Architecture}

The DOI system integrates with the EBF SuperKey architecture:

\begin{table}[h]
\centering
\caption{SuperKey Prefixes}
\begin{tabular}{lll}
\toprule
\textbf{Prefix} & \textbf{Entity} & \textbf{Format} \\
\midrule
PAP- & Paper & \texttt{PAP-\{author\}\{year\}\{keyword\}} \\
JRN- & Journal & \texttt{JRN-\{registrant\}-\{abbrev\}} \\
DOI- & DOI & \texttt{10.\{registrant\}/\{suffix\}} \\
\bottomrule
\end{tabular}
\end{table}

\subsection{Workflow Integration}

\begin{enumerate}
    \item \textbf{New Paper Entry:} Run \texttt{doi\_strategy.py --key \{key\}} to determine strategy
    \item \textbf{PATTERN Result:} Add generated DOI to \texttt{bcm\_master.bib}
    \item \textbf{LOOKUP Result:} Query CrossRef/JSTOR, then add DOI
    \item \textbf{Validation:} Verify DOI resolves correctly via \texttt{doi.org}
\end{enumerate}

% ============================================================================
\section{Limitations and Future Work}
\label{sec:limitations}
% ============================================================================

\subsection{Current Limitations}

\begin{enumerate}
    \item \textbf{Partial Coverage:} Only 45 of $\sim$100+ economics journals modeled
    \item \textbf{Era Boundaries:} Transition years are approximate ($\pm$2 years)
    \item \textbf{Pattern Validation:} Generated DOIs should be validated against CrossRef
    \item \textbf{API Blocking:} CrossRef/JSTOR APIs may rate-limit bulk requests
\end{enumerate}

\subsection{Future Extensions}

\begin{enumerate}
    \item Expand journal database to cover all journals in bibliography
    \item Implement automated CrossRef validation for generated DOIs
    \item Add ISBN lookup for books (currently SKIPped)
    \item Integrate with Zotero/Mendeley for automated bibliography management
\end{enumerate}

% ============================================================================
\section{Conclusion}
\label{sec:conclusion}
% ============================================================================

The DOI prediction problem can be decomposed into a journal-era model where predictability depends on:

\begin{enumerate}
    \item The journal's publisher and their suffix algorithm
    \item The publication year and applicable era
    \item Availability of required metadata (volume, issue, pages)
\end{enumerate}

For the EBF bibliography:
\begin{itemize}
    \item \textbf{24\%} of journals have fully predictable DOIs (forecast error < 10\%)
    \item \textbf{38\%} have partially predictable DOIs (era-dependent)
    \item \textbf{38\%} require external lookup (hash-based or legacy systems)
\end{itemize}

The \texttt{doi\_strategy.py} script operationalizes this model, enabling systematic DOI completion for the entire bibliography.

% ============================================================================
\section{Summary}
\label{sec:summary}
% ============================================================================

This appendix presented a systematic approach to DOI management in the EBF bibliography:

\begin{enumerate}
    \item \textbf{ISO 26324 Standard:} DOIs follow the structure \texttt{10.RRRR/SUFFIX}, where the prefix is standardized but the suffix is publisher-specific.

    \item \textbf{Publisher Algorithms:} Six pattern types were identified (STRUCTURED, ISSN\_YEAR, NUMERIC, LONG\_CODE, ALPHA\_NUM, OTHER), with predictability varying by publisher and era.

    \item \textbf{Journal-Era Model:} Each journal has distinct eras with different DOI patterns. The predictability rate $R_j$ quantifies the proportion of papers with deterministic DOI generation.

    \item \textbf{Strategy Determination:} The \texttt{doi\_strategy.py} script implements Algorithm~\ref{alg:doi-strategy} to classify each paper as PATTERN (generate) or LOOKUP (query API).

    \item \textbf{Results:} For 1,604 papers without DOI:
    \begin{itemize}
        \item 11\% can have DOIs generated from metadata
        \item 72\% require API lookup (CrossRef, JSTOR)
        \item 17\% are books/proceedings (skipped)
    \end{itemize}
\end{enumerate}

The journal database (\texttt{journal-database.yaml} v3.0) provides the foundation for systematic DOI completion.

% ============================================================================
\section{Glossary}
\label{sec:glossary}
% ============================================================================

Key terms used in this appendix (see also Appendix G for comprehensive glossary):

\begin{description}
    \item[DOI] Digital Object Identifier -- a persistent identifier for digital objects, standardized in ISO 26324.

    \item[Registrant] An organization authorized to assign DOIs within a specific prefix (e.g., 10.1257 = American Economic Association).

    \item[Registration Agency (RA)] An organization that assigns DOI prefixes to registrants (e.g., CrossRef for scholarly publications).

    \item[Suffix] The publisher-specific part of a DOI following the prefix and slash.

    \item[Predictability Rate] The proportion of papers in a journal for which DOIs can be generated from metadata.

    \item[Forecast Error] $1 - R_j$; the proportion of papers requiring API lookup.

    \item[Era] A time period during which a journal used a consistent DOI pattern.

    \item[SuperKey] EBF's unique identifier format: JRN- for journals, PAP- for papers.

    \item[JSTOR] A digital library providing archival access to academic journals; assigns DOIs with prefix 10.2307.

    \item[CrossRef] The primary DOI registration agency for scholarly publications.
\end{description}

% ============================================================================
\section{Critical Foundations}
\label{sec:foundations}
% ============================================================================

The DOI reference system rests on three foundations:

\begin{enumerate}
    \item \textbf{ISO 26324 Compliance:} All DOIs in the EBF bibliography conform to the international standard syntax.

    \item \textbf{Empirical Pattern Analysis:} Publisher algorithms were reverse-engineered from 720 known DOIs in the bibliography, not assumed.

    \item \textbf{Conservative Prediction:} When uncertain, the system defaults to LOOKUP rather than generating potentially invalid DOIs.
\end{enumerate}

% ============================================================================
\section{Open Issues}
\label{sec:open-issues}
% ============================================================================

\begin{enumerate}
    \item \textbf{OI-1: Expanded Journal Coverage} -- Currently 45 of $\sim$100+ journals are modeled. Priority: High.

    \item \textbf{OI-2: API Integration} -- Automated CrossRef/JSTOR lookup not yet implemented. Priority: Medium.

    \item \textbf{OI-3: Validation Pipeline} -- Generated DOIs should be validated against CrossRef before storage. Priority: High.

    \item \textbf{OI-4: Book ISBN Support} -- 279 book entries currently skipped; ISBN lookup needed. Priority: Low.
\end{enumerate}

% ----------------------------------------------------------------------------
% REFERENCES
% ----------------------------------------------------------------------------
\section{References}
\label{sec:references}

\begin{itemize}
    \item ISO 26324:2022. \textit{Information and documentation -- Digital object identifier system}. International Organization for Standardization.
    \item CrossRef. \textit{Constructing your DOIs}. \url{https://www.crossref.org/documentation/member-setup/constructing-your-dois/}
    \item DataCite. \textit{DOI Basics}. \url{https://support.datacite.org/docs/doi-basics}
    \item International DOI Foundation. \textit{DOI Handbook}. \url{https://www.doi.org/doi-handbook/}
\end{itemize}

% Link to master bibliography
\nocite{bcm_master}

\end{chapter}
